% ==========================================================================
% Midterm report - Supervised Learning section (drop-in for Overleaf).
% Copy supervised/outputs/figures/*.png into the Overleaf project first:
%   ridge_pred_vs_true.png, logistic_pr_curve.png, logistic_roc_curve.png
% ==========================================================================

\subsection{Supervised Learning: Baseline Prediction Models}

\textbf{Setup.} For the first supervised milestone, we implement two
deliberately simple baselines from the proposal: Ridge regression for
predicting 3-year citation counts and Logistic Regression for identifying
highly cited papers. This section establishes a leakage-safe baseline before
adding the cluster-derived features suggested by the unsupervised analysis.
Both models therefore reuse the same base representation as the
unsupervised-stage wherever possible. Title + abstract text is featurized
with TF-IDF, compressed to 100 LSA dimensions with randomized Truncated SVD,
and L2-normalized; metadata features include author count, reference count,
venue frequency, year, category indicators, affiliation/country counts, and
title/abstract lengths. The key supervised change is that all preprocessing
is fit only on the c1 train split (24{,}691 papers) and then applied unchanged
to the held-out temporal test split (25{,}308 papers), avoiding leakage from
future papers.

\textbf{Model settings.} Ridge uses $\alpha{=}1.0$ and is trained on
$\log(1+\text{citations})$. Logistic Regression uses $C{=}1.0$,
\texttt{class\_weight="balanced"}, and \texttt{max\_iter=2000}. Class
balancing is used because the high-citation class is rare under the
train-derived top-decile threshold. For the reported F1, precision, recall,
and accuracy, Logistic Regression probabilities are converted to labels with
the default threshold $p \ge 0.5$.

\textbf{Regression baseline.} Ridge is trained on
$\log(1+\text{3-year citations})$. On the test split it reaches RMSE =
1.111, MAE = 0.880, $R^2$ = 0.141, and Spearman rank correlation = 0.395
(Table~\ref{tab:ridge-regression}). The modest $R^2$ shows that exact
citation counts are hard to predict with a linear baseline, especially
because citation counts are heavy-tailed. However, the positive Spearman
correlation indicates that the model is useful as a ranking signal.

\begin{table}[t]
\centering
\begin{tabular}{llrrrr}
\hline
Model & Split & RMSE & MAE & $R^2$ & Spearman \\
\hline
Ridge & Train & 1.193 & 0.956 & 0.179 & 0.411 \\
Ridge & Test  & 1.111 & 0.880 & 0.141 & 0.395 \\
\hline
\end{tabular}
\caption{Ridge regression performance on $\log(1+\text{citations})$.}
\label{tab:ridge-regression}
\end{table}

\begin{figure}[t]
\centering
\includegraphics[width=0.82\linewidth]{ridge_pred_vs_true.png}
\caption{Ridge predicted versus true $\log(1+\text{citations})$ on the test
split. The model captures broad ordering but compresses predictions toward
the mean, which is expected for an L2-regularized linear model on
heavy-tailed citation data.}
\label{fig:ridge-pred-true}
\end{figure}

\textbf{Ranking evaluation.} The proposal emphasizes ranking metrics because
the practical goal is often to surface papers likely to become influential
rather than to predict the exact citation count. Using the Ridge score to
rank test papers, the top 10\% contains highly cited papers at 20.0\%
precision, compared with a 5.9\% base rate in the test split
(Table~\ref{tab:ridge-ranking}). This corresponds to a 3.39x lift.
Recall@10\% is 33.9\%, and NDCG@10\% is 0.633. These results are consistent
with the unsupervised finding that topic and metadata contain meaningful
impact signal even before cluster id or distance-to-centroid features are
added.

\begin{table}[t]
\centering
\begin{tabular}{llrrrr}
\hline
Model & Split & Precision@10\% & Recall@10\% & NDCG@10\% & Lift@10\% \\
\hline
Ridge & Train & 0.297 & 0.294 & 0.655 & 2.94x \\
Ridge & Test  & 0.200 & 0.339 & 0.633 & 3.39x \\
\hline
\end{tabular}
\caption{Ranking performance when Ridge predictions are used to rank papers.}
\label{tab:ridge-ranking}
\end{table}

\textbf{Classification baseline.} We define ``highly cited'' using the train
split top-decile threshold: at least 41 citations within 3 years. The
temporal test set is more difficult, with only 5.9\% positives under this
fixed threshold. On test, Logistic Regression obtains ROC-AUC = 0.773 and
PR-AUC = 0.180, compared with a PR baseline of 0.059 from prevalence alone.
At the default threshold the model has F1 = 0.255, precision = 0.169, and
recall = 0.523 (Table~\ref{tab:logistic-classification}). This is a
reasonable first-stage screening model: it finds over half of the highly
cited test papers, but with many false positives.

\begin{table}[t]
\centering
\begin{tabular}{llrrrrr}
\hline
Model & Split & PR-AUC & ROC-AUC & F1 & Precision & Recall \\
\hline
Logistic Regression & Train & 0.263 & 0.771 & 0.316 & 0.202 & 0.719 \\
Logistic Regression & Test  & 0.180 & 0.773 & 0.255 & 0.169 & 0.523 \\
Prevalence baseline & Test  & 0.059 & 0.500 & 0.000 & 0.000 & 0.000 \\
\hline
\end{tabular}
\caption{Binary classification performance for predicting high-citation
papers using a train-split threshold of 41 citations.}
\label{tab:logistic-classification}
\end{table}

\begin{figure}[t]
\centering
\includegraphics[width=0.74\linewidth]{logistic_pr_curve.png}
\caption{Precision-recall curve for the Logistic Regression baseline. The
PR-AUC of 0.180 is about 3.0 times the 0.059 positive base rate in the test
set.}
\label{fig:logistic-pr}
\end{figure}

\begin{figure}[t]
\centering
\includegraphics[width=0.74\linewidth]{logistic_roc_curve.png}
\caption{ROC curve for the Logistic Regression baseline. ROC-AUC is 0.773,
indicating useful discrimination even though the positive class is rare.}
\label{fig:logistic-roc}
\end{figure}

\textbf{Cluster-feature ablation.} To test the direct link between the
unsupervised and supervised stages, we export the text-only $k{=}8$ K-Means
model from the unsupervised pipeline and append two feature groups to the
supervised matrices: cluster-id one-hot indicators and distances to all eight
centroids. These features are fit only on the train split by the unsupervised
stage and then applied to the test split. The augmented matrix has 260
features instead of 244. On the test split, however, the changes are
negligible: Ridge RMSE changes from 1.111 to 1.111, Spearman from 0.3951 to
0.3951, Precision@10\% from 0.2003 to 0.2007, and Logistic PR-AUC from 0.1798
to 0.1794. This suggests that the linear baselines already absorb most of the
topic information through the continuous LSA features. The clusters are still
valuable for interpretation and per-cluster error analysis, but they do not
materially strengthen these first-stage linear models as additional features.

\textbf{Next steps.} These baselines establish a leakage-safe supervised
pipeline and cover the proposal's regression, ranking, and
binary-classification metrics. The main weakness is calibration of the
citation tail: Ridge compresses high-impact papers toward average
predictions. In the next stage, we will add stronger nonlinear models and
report per-cluster error analysis so that model performance can be compared
across impact archetypes rather than only globally.
